\ifdefined\idpCombinedMode
\else
\documentclass{idpbeamer}
\fi


\lectureinfo{POO Java}{Aula 01\\Programação Orientada a Objetos em Java}{Prof. Dr. Nome do Professor}

\startLectureFile

\providecommand{\pooCodePath}{../code}

\section{Objetivos}

\begin{frame}{Objetivos}
Nesta aula, vamos organizar o vocabulário inicial da disciplina e preparar o ambiente para os primeiros programas em Java.

\begin{itemize}
  \item Entender o papel de classes, objetos e métodos.
  \item Compilar e executar um programa Java simples.
  \item Identificar a estrutura mínima de uma aplicação.
\end{itemize}
\end{frame}

\section{Primeiro Programa}

\begin{frame}[fragile]{Primeiro Programa}
\lstinputlisting{\pooCodePath/OlaMundo.java}

O nome do arquivo deve acompanhar o nome da classe pública: \texttt{OlaMundo.java}.
\end{frame}

\section{Próximos Passos}

\begin{frame}{Próximos Passos}
Ao final da aula, registre dúvidas e exemplos que merecem voltar na próxima discussão.

\begin{block}{Para praticar}
Crie uma classe com seu nome e faça o programa imprimir uma pequena apresentação no terminal.
\end{block}
\end{frame}

\finishLectureFile
